\documentclass{report}

\usepackage{graphicx}
\usepackage{algorithm}
\usepackage{algpseudocode}
\usepackage{float}
\usepackage{adjustbox}
\graphicspath{{graphics/}}

\begin{document}
	\begin{titlepage}
		\centering
		\includegraphics[scale=0.35]{logo_feup.png}\linebreak
		
		\vspace{1cm}
		
		{\scshape \large Bachelor in Informatics and Computing Engineering}
		
		\vspace {1cm}
		
		{\scshape\Huge Distributed and Partitioned Key-Value Store \par}
		
		\vfill
		
		{\scshape \large Parallel and Distributed Computing}
		
		\vfill
		
		\Large David \textsc{Preda} - up201904726 \\ Fernando
		\textsc{Rego} - up201905951 \\ Miguel \textsc{Amorim} - up201907756
		
		\vspace{1cm}
		
		\today
		
	\end{titlepage}

	\tableofcontents
	
	\chapter{Problem Description}
			\paragraph{} A key-value store is a simple storage system that stores
			arbitrary data objects, the values, each of which is accessed by means
			of a key, very much like in a hash table. To ensure persistency, the data
			items and their keys must be stored in persistent storage, e.g. a hard
			disk drive (HDD) or a solid state disk (SSD), rather than in RAM.
			
			\paragraph{}By distributed, we mean that the data items in the key-value
			store are partitioned among different cluster nodes.
			
			\paragraph{} Our design is loosely based on Amazon's Dynamo, in that it
			uses consistent-hashing to partition the key-value pairs among the 
			different nodes.
	
	\chapter{Message Format}

			\paragraph{}
	        
	        \paragraph{}

	\chapter{Membership Service}
	
	        \paragraph{}
	        
	        \paragraph{}

			\section{Implementation}
	
				\paragraph{}
				
				\paragraph{}
	
	\chapter{Key-value Store}
	
	        \paragraph{}
	        
	        \paragraph{}

			\section{Implementation}
	
				\paragraph{}
				
				\paragraph{}
	
	\chapter{Replication}
	
	        \paragraph{}
	        
	        \paragraph{}

			\section{Implementation}
	
				\paragraph{}
				
				\paragraph{}

			\section{Implications on membership and storage devices}
	
				\paragraph{}
				
				\paragraph{}

	\chapter{Fault-Tolerance}
	
	        \paragraph{}
	        
	        \paragraph{}

	\chapter{Thread-pools}
	
	        \paragraph{}
	        
	        \paragraph{}

			\section{Implementation}
	
				\paragraph{}
				
				\paragraph{}

	\chapter{Test Client}
	
	        \paragraph{}
	        
	        \paragraph{}

	\chapter{Conclusions}
	
	        \paragraph{}This project allowed us to deeply understand some computer concepts related to store .
	        
	        \paragraph{}
			

\end{document}
